% Mathematics Bootcamp --- Day 1 Exercises
% Mas-Colell, Whinston and Green, Chapter 1, Sections 1.B and 1.C.
%
% Self-contained: this file needs no other file in the repository.
% Compile with pdflatex, twice, or with latexmk -pdf.

\documentclass[12pt]{article}

\usepackage[margin=1in]{geometry}

\usepackage{
  amsmath,
  amssymb,
  amsfonts,
  amsthm,
}

\usepackage[T1]{fontenc}
\usepackage{lmodern}

\usepackage{setspace}
\usepackage{float}

% Added to the template: the roman/alpha labels in the exercise statements
% need it. source/proofs/main.tex already loads it too.
\usepackage{enumitem}

\onehalfspacing

% ---------------------------------------------------------------------------
% The only additions to the template: a heading for each exercise, a
% "Need to Show" block, and one level of indentation for a structured proof.
% ---------------------------------------------------------------------------

% \exercise{number}{MWG grade}{one-line description}
\newcommand{\exercise}[3]{%
  \par\penalty-250\bigskip
  \noindent\textbf{\large Exercise #1}\quad\textnormal{(grade #2)}%
  \par\nobreak\smallskip
  \noindent\textit{#3}%
  \par\nobreak\smallskip
  \noindent\rule{\linewidth}{0.4pt}%
  \par\nobreak\smallskip}

% The statement, verbatim from the book.
\newenvironment{statement}{\par\medskip}{\par\medskip}

% What the proof is obliged to produce, written before any of it is written.
\newenvironment{nts}
  {\par\medskip\noindent\textbf{Need to Show:}\par\nobreak\smallskip}
  {\par\medskip}

% One level of indentation. Nest to indent further: the structure of the
% argument is carried by the indentation, as in the Day 1 slides.
\newlength{\plevel}\setlength{\plevel}{2.25em}
\newenvironment{pstep}
  {\begin{list}{}{%
     \setlength{\leftmargin}{\plevel}%
     \setlength{\rightmargin}{0pt}%
     \setlength{\labelwidth}{0pt}\setlength{\labelsep}{0pt}%
     \setlength{\itemindent}{0pt}\setlength{\listparindent}{0pt}%
     \setlength{\topsep}{0pt}\setlength{\partopsep}{0pt}%
     \setlength{\parsep}{0pt}\setlength{\itemsep}{0pt}}%
   \item\relax}
  {\end{list}}

% Right-flushed justification for a line of a proof. Ends the line.
\newcommand{\by}[1]{%
  \unskip\nobreak\hfil\penalty50\hskip1em\null\nobreak\hfil
  \textnormal{[#1]}%
  {\parfillskip=0pt \finalhyphendemerits=0 \par}}

\newenvironment{remark}
  {\par\medskip\noindent\textit{Remark.}\enspace}
  {\par\medskip}

% Contradiction. The single arrow is the one the Day 1 slides use for
% implication, so this reads as "then, and then back the other way".
\newcommand{\contra}{\ensuremath{\rightarrow\leftarrow}}

\title{\textbf{Mathematics Bootcamp --- Day 1 Exercises}\\[.4em]
       \large Mas-Colell, Whinston and Green, Chapter 1}
\author{Miguel V\'azquez-Carrero}
\date{Preference and Choice: Sections 1.B and 1.C\\[.3em]
      \normalsize Compilation Date: \today}

\begin{document}

\maketitle

\section*{Notation}

Throughout, $X$ is a set of alternatives and $\succeq$ is a binary relation on
$X$, read ``at least as good as''. From it we derive
\[
x\succ y \iff x\succeq y \ \text{ and not } y\succeq x,
\qquad
x\sim y \iff x\succeq y \ \text{ and } y\succeq x .
\]
The relation $\succeq$ is \emph{rational} if it is
\emph{complete} (for all $x,y\in X$, $x\succeq y$ or $y\succeq x$)
and \emph{transitive} (for all $x,y,z\in X$, if $x\succeq y$ and $y\succeq z$
then $x\succeq z$). A function $u:X\to\mathbb R$ \emph{represents} $\succeq$ if
\[
x\succeq y \iff u(x)\geq u(y)\qquad\text{for all }x,y\in X .
\]

A \emph{choice structure} is a pair $(\mathbb B,C(\cdot))$: $\mathbb B$ is a
family of nonempty subsets of $X$, the observed menus, and $C$ assigns to each
$B\in\mathbb B$ a nonempty set $C(B)\subseteq B$. MWG writes $\mathcal B$ for
the menu domain; we keep the $\mathbb B$ of the slides.

\begin{quote}
\textbf{Definition 1.C.1 (Weak Axiom).}
If for some $B\in\mathbb B$ with $x,y\in B$ we have $x\in C(B)$, then for any
$B^{\prime}\in\mathbb B$ with $x,y\in B^{\prime}$ and $y\in C(B^{\prime})$, we
must also have $x\in C(B^{\prime})$.

\medskip
\textbf{Definition 1.C.2 (Revealed preference).}
$x\succeq^{*}y \iff$ there is some $B\in\mathbb B$ such that $x,y\in B$ and
$x\in C(B)$.
\end{quote}

Exercises 1.B.1 and 1.B.2 ask for the proof of the following, which MWG states
and does not prove.

\begin{quote}
\textbf{Proposition 1.B.1.} If $\succeq$ is rational then:
\begin{enumerate}[label=(\roman*),leftmargin=2.4em,topsep=.3em,itemsep=.2em]
  \item $\succ$ is both \emph{irreflexive} ($x\succ x$ never holds) and
        \emph{transitive} (if $x\succ y$ and $y\succ z$, then $x\succ z$).
  \item $\sim$ is \emph{reflexive} ($x\sim x$ for all $x$), \emph{transitive}
        (if $x\sim y$ and $y\sim z$, then $x\sim z$), and \emph{symmetric}
        (if $x\sim y$, then $y\sim x$).
  \item if $x\succ y\succeq z$, then $x\succ z$.
\end{enumerate}
\end{quote}

\section*{How to write these up}

Write every proof the way a program is written: the skeleton first, then the
body. Recall the skeleton from Day 1,
\begin{quote}
Let $x$ be arbitrary.\par
\hspace*{2.25em}Suppose $P(x)$ is true.\par
\hspace*{4.5em}[Proof of $Q(x)$ goes here.]\par
\hspace*{2.25em}Thus, if $P(x)$ then $Q(x)$.\par
Thus, for all $x$, if $P(x)$ then $Q(x)$.
\end{quote}
and note that the indentation is not decoration. A line is available exactly to
the lines indented under it, and a block is closed by the line that returns to
the outer margin. Write the outer lines first: they are forced by the shape of
the statement alone, before you have any idea why it is true.

Before any of that, state what the proof owes. Unfold every definition in the
conclusion until what is left is a list of claims with no defined terms in
them. That list is the \textbf{Need to Show}. A proof is finished when the list
is exhausted, and not before.

\newpage
\section*{Preferences (Section 1.B)}

% ===========================================================================
\exercise{1.B.1}{B}{Property (iii) of Proposition 1.B.1}

\begin{statement}
Prove property (iii) of Proposition 1.B.1.
\end{statement}


% ===========================================================================
\exercise{1.B.2}{A}{Properties (i) and (ii) of Proposition 1.B.1}

\begin{statement}
Prove properties (i) and (ii) of Proposition 1.B.1.
\end{statement}


% ===========================================================================
\exercise{1.B.3}{B}{A strictly increasing transformation of a utility function}

\begin{statement}
Show that if $f:\mathbb R\to\mathbb R$ is a strictly increasing function and
$u:X\to\mathbb R$ is a utility function representing preference relation
$\succeq$, then the function $v:X\to\mathbb R$ defined by $v(x)=f(u(x))$ is
also a utility function representing preference relation $\succeq$.
\end{statement}


% ===========================================================================
\exercise{1.B.4}{A}{A sufficient condition for representation}

\begin{statement}
Consider a rational preference relation $\succeq$. Show that if $u(x)=u(y)$
implies $x\sim y$ and if $u(x)>u(y)$ implies $x\succ y$, then $u(\cdot)$ is a
utility function representing $\succeq$.
\end{statement}


% ===========================================================================
\exercise{1.B.5}{B}{Representation on a finite set}

\begin{statement}
Show that if $X$ is finite and $\succeq$ is a rational preference relation on
$X$, then there is a utility function $u:X\to\mathbb R$ that represents
$\succeq$. [\emph{Hint}: Consider first the case in which the individual's
ranking between any two elements of $X$ is strict (i.e., there is never any
indifference), and construct a utility function representing these preferences;
then extend your argument to the general case.]
\end{statement}


\newpage
\section*{Choice Rules (Section 1.C)}

% ===========================================================================
\exercise{1.C.1}{B}{What the weak axiom leaves open}

\begin{statement}
Consider the choice structure $(\mathbb B,C(\cdot))$ with
$\mathbb B=\bigl\{\{x,y\},\{x,y,z\}\bigr\}$ and $C(\{x,y\})=\{x\}$. Show that
if $(\mathbb B,C(\cdot))$ satisfies the weak axiom, then we must have
$C(\{x,y,z\})=\{x\}$, $=\{z\}$, or $=\{x,z\}$.
\end{statement}


% ===========================================================================
\exercise{1.C.2}{B}{A symmetric restatement of the weak axiom}

\begin{statement}
Show that the weak axiom (Definition 1.C.1) is equivalent to the following
property holding:
\begin{quote}
Suppose that $B,B^{\prime}\in\mathbb B$, that $x,y\in B$, and that
$x,y\in B^{\prime}$. Then if $x\in C(B)$ and $y\in C(B^{\prime})$, we must have
$\{x,y\}\subseteq C(B)$ and $\{x,y\}\subseteq C(B^{\prime})$.
\end{quote}
\end{statement}


% ===========================================================================
\exercise{1.C.3}{C}{Two revealed strict preferences}

\begin{statement}
Suppose that choice structure $(\mathbb B,C(\cdot))$ satisfies the weak axiom.
Consider the following two possible revealed preferred relations, $\succ^{*}$
and $\succ^{**}$:
\[
\begin{aligned}
x\succ^{*}y \quad&\Longleftrightarrow\quad
\text{there is some }B\in\mathbb B\text{ such that }x,y\in B,\ x\in C(B),
\text{ and } y\notin C(B)\\
x\succ^{**}y \quad&\Longleftrightarrow\quad
x\succeq^{*}y\text{ but not }y\succeq^{*}x
\end{aligned}
\]
where $\succeq^{*}$ is the revealed at-least-as-good-as relation defined in
Definition 1.C.2.
\begin{enumerate}[label=(\alph*),leftmargin=2.4em,topsep=.3em,itemsep=.25em]
  \item Show that $\succ^{*}$ and $\succ^{**}$ give the same relation over $X$;
        that is, for any $x,y\in X$, $x\succ^{*}y\iff x\succ^{**}y$. Is this
        still true if $(\mathbb B,C(\cdot))$ does not satisfy the weak axiom?
  \item Must $\succ^{*}$ be transitive?
  \item Show that if $\mathbb B$ includes all three-element subsets of $X$,
        then $\succ^{*}$ is transitive.
\end{enumerate}
\end{statement}

\end{document}
